% Options for packages loaded elsewhere
\PassOptionsToPackage{unicode}{hyperref}
\PassOptionsToPackage{hyphens}{url}
\PassOptionsToPackage{dvipsnames,svgnames,x11names}{xcolor}
%
\documentclass[
  10pt,
]{article}
\usepackage{amsmath,amssymb}
\usepackage{iftex}
\ifPDFTeX
  \usepackage[T1]{fontenc}
  \usepackage[utf8]{inputenc}
  \usepackage{textcomp} % provide euro and other symbols
\else % if luatex or xetex
  \usepackage{unicode-math} % this also loads fontspec
  \defaultfontfeatures{Scale=MatchLowercase}
  \defaultfontfeatures[\rmfamily]{Ligatures=TeX,Scale=1}
\fi
\usepackage{lmodern}
\ifPDFTeX\else
  % xetex/luatex font selection
\fi
% Use upquote if available, for straight quotes in verbatim environments
\IfFileExists{upquote.sty}{\usepackage{upquote}}{}
\IfFileExists{microtype.sty}{% use microtype if available
  \usepackage[]{microtype}
  \UseMicrotypeSet[protrusion]{basicmath} % disable protrusion for tt fonts
}{}
\makeatletter
\@ifundefined{KOMAClassName}{% if non-KOMA class
  \IfFileExists{parskip.sty}{%
    \usepackage{parskip}
  }{% else
    \setlength{\parindent}{0pt}
    \setlength{\parskip}{6pt plus 2pt minus 1pt}}
}{% if KOMA class
  \KOMAoptions{parskip=half}}
\makeatother
\usepackage{xcolor}
\usepackage[margin=0.82in]{geometry}
\usepackage{color}
\usepackage{fancyvrb}
\newcommand{\VerbBar}{|}
\newcommand{\VERB}{\Verb[commandchars=\\\{\}]}
\DefineVerbatimEnvironment{Highlighting}{Verbatim}{commandchars=\\\{\}}
% Add ',fontsize=\small' for more characters per line
\newenvironment{Shaded}{}{}
\newcommand{\AlertTok}[1]{\textcolor[rgb]{1.00,0.00,0.00}{\textbf{#1}}}
\newcommand{\AnnotationTok}[1]{\textcolor[rgb]{0.38,0.63,0.69}{\textbf{\textit{#1}}}}
\newcommand{\AttributeTok}[1]{\textcolor[rgb]{0.49,0.56,0.16}{#1}}
\newcommand{\BaseNTok}[1]{\textcolor[rgb]{0.25,0.63,0.44}{#1}}
\newcommand{\BuiltInTok}[1]{\textcolor[rgb]{0.00,0.50,0.00}{#1}}
\newcommand{\CharTok}[1]{\textcolor[rgb]{0.25,0.44,0.63}{#1}}
\newcommand{\CommentTok}[1]{\textcolor[rgb]{0.38,0.63,0.69}{\textit{#1}}}
\newcommand{\CommentVarTok}[1]{\textcolor[rgb]{0.38,0.63,0.69}{\textbf{\textit{#1}}}}
\newcommand{\ConstantTok}[1]{\textcolor[rgb]{0.53,0.00,0.00}{#1}}
\newcommand{\ControlFlowTok}[1]{\textcolor[rgb]{0.00,0.44,0.13}{\textbf{#1}}}
\newcommand{\DataTypeTok}[1]{\textcolor[rgb]{0.56,0.13,0.00}{#1}}
\newcommand{\DecValTok}[1]{\textcolor[rgb]{0.25,0.63,0.44}{#1}}
\newcommand{\DocumentationTok}[1]{\textcolor[rgb]{0.73,0.13,0.13}{\textit{#1}}}
\newcommand{\ErrorTok}[1]{\textcolor[rgb]{1.00,0.00,0.00}{\textbf{#1}}}
\newcommand{\ExtensionTok}[1]{#1}
\newcommand{\FloatTok}[1]{\textcolor[rgb]{0.25,0.63,0.44}{#1}}
\newcommand{\FunctionTok}[1]{\textcolor[rgb]{0.02,0.16,0.49}{#1}}
\newcommand{\ImportTok}[1]{\textcolor[rgb]{0.00,0.50,0.00}{\textbf{#1}}}
\newcommand{\InformationTok}[1]{\textcolor[rgb]{0.38,0.63,0.69}{\textbf{\textit{#1}}}}
\newcommand{\KeywordTok}[1]{\textcolor[rgb]{0.00,0.44,0.13}{\textbf{#1}}}
\newcommand{\NormalTok}[1]{#1}
\newcommand{\OperatorTok}[1]{\textcolor[rgb]{0.40,0.40,0.40}{#1}}
\newcommand{\OtherTok}[1]{\textcolor[rgb]{0.00,0.44,0.13}{#1}}
\newcommand{\PreprocessorTok}[1]{\textcolor[rgb]{0.74,0.48,0.00}{#1}}
\newcommand{\RegionMarkerTok}[1]{#1}
\newcommand{\SpecialCharTok}[1]{\textcolor[rgb]{0.25,0.44,0.63}{#1}}
\newcommand{\SpecialStringTok}[1]{\textcolor[rgb]{0.73,0.40,0.53}{#1}}
\newcommand{\StringTok}[1]{\textcolor[rgb]{0.25,0.44,0.63}{#1}}
\newcommand{\VariableTok}[1]{\textcolor[rgb]{0.10,0.09,0.49}{#1}}
\newcommand{\VerbatimStringTok}[1]{\textcolor[rgb]{0.25,0.44,0.63}{#1}}
\newcommand{\WarningTok}[1]{\textcolor[rgb]{0.38,0.63,0.69}{\textbf{\textit{#1}}}}
\usepackage{longtable,booktabs,array}
\usepackage{calc} % for calculating minipage widths
% Correct order of tables after \paragraph or \subparagraph
\usepackage{etoolbox}
\makeatletter
\patchcmd\longtable{\par}{\if@noskipsec\mbox{}\fi\par}{}{}
\makeatother
% Allow footnotes in longtable head/foot
\IfFileExists{footnotehyper.sty}{\usepackage{footnotehyper}}{\usepackage{footnote}}
\makesavenoteenv{longtable}
\usepackage{graphicx}
\makeatletter
\def\maxwidth{\ifdim\Gin@nat@width>\linewidth\linewidth\else\Gin@nat@width\fi}
\def\maxheight{\ifdim\Gin@nat@height>\textheight\textheight\else\Gin@nat@height\fi}
\makeatother
% Scale images if necessary, so that they will not overflow the page
% margins by default, and it is still possible to overwrite the defaults
% using explicit options in \includegraphics[width, height, ...]{}
\setkeys{Gin}{width=\maxwidth,height=\maxheight,keepaspectratio}
% Set default figure placement to htbp
\makeatletter
\def\fps@figure{htbp}
\makeatother
\setlength{\emergencystretch}{3em} % prevent overfull lines
\providecommand{\tightlist}{%
  \setlength{\itemsep}{0pt}\setlength{\parskip}{0pt}}
\setcounter{secnumdepth}{-\maxdimen} % remove section numbering
\ifLuaTeX
\usepackage[bidi=basic]{babel}
\else
\usepackage[bidi=default]{babel}
\fi
\babelprovide[main,import]{american}
% get rid of language-specific shorthands (see #6817):
\let\LanguageShortHands\languageshorthands
\def\languageshorthands#1{}
\usepackage{xurl}
\usepackage{microtype}
\usepackage{fancyhdr}
\pagestyle{fancy}
\fancyhf{}
\fancyhead[L]{\small ProofNet / Joint research design}
\fancyhead[R]{\small 2026-09-30}
\fancyfoot[C]{\thepage}
\setlength{\headheight}{15pt}
\emergencystretch=3em
\ifLuaTeX
  \usepackage{selnolig}  % disable illegal ligatures
\fi
\usepackage{bookmark}
\IfFileExists{xurl.sty}{\usepackage{xurl}}{} % add URL line breaks if available
\urlstyle{same}
\hypersetup{
  pdftitle={From Proof-Order Graphs to Checked Proof Reuse},
  pdfauthor={Research design note for the ProofNet collaboration},
  pdflang={en-US},
  colorlinks=true,
  linkcolor={black},
  filecolor={Maroon},
  citecolor={Blue},
  urlcolor={blue},
  pdfcreator={LaTeX via pandoc}}

\title{From Proof-Order Graphs to Checked Proof Reuse}
\usepackage{etoolbox}
\makeatletter
\providecommand{\subtitle}[1]{% add subtitle to \maketitle
  \apptocmd{\@title}{\par {\large #1 \par}}{}{}
}
\makeatother
\subtitle{A source-pinned comparison and joint research plan for
ProofNet}
\author{Research design note for the ProofNet collaboration}
\date{30 September 2026}

\begin{document}
\maketitle
\begin{abstract}
This note compares the current ProofNet-IR and proof-graphs repositories
with the available Boundary-Lens prototype. The repositories already
supply a verified restricted MLL library, Lean graph extraction,
structural identities, coupled-goal search, and negative empirical
results that narrow the case for further order-based pruning. We
recommend one integrated program: retain the formal MLL layer, extend
the existing Lean engine with typed transaction-aware boundaries, and
test completed proof-fragment transport against a conventional typed
cache. We separate standard substitution principles from proposed
engineering and distinguish final kernel replay from observed tactic
effects. New executable results here are a Python audit of the supplied
prototype and finite-model regression tests, not fresh Lean or
neural-prover experiments. Source pins, explicit evidence levels,
falsifiable experiments, and implementation prompts make the proposed
collaboration auditable.
\end{abstract}

{
\hypersetup{linkcolor=}
\setcounter{tocdepth}{2}
\tableofcontents
}
\section{Reading this note}\label{reading-this-note}

The conclusions compare inspected artifacts, not the collaborators'
personal abilities or complete contribution histories. Numerical
repository results below are first-party reports, not new replications.
The review could read selected current GitHub files, but the container
had no Lean/Lake and could not clone the repositories. The local
user-artifact audit and the reference Python tests were executed. All
new live-graph and proof-reuse examples are design fixtures.

The package contains 29 passing Python representation/schema tests.
Those tests exercise a deliberately restricted expression model. They
neither typecheck Lean terms nor prove that a proposed normalization
preserves all of Lean's operational behavior.

\begin{figure}
\centering
\includegraphics[width=0.88\textwidth,height=\textheight]{../figures/joint_architecture.pdf}
\caption{One integration, not three competing extractors. The MLL
reference arrows express reuse of test cases and concepts, not a
soundness reduction for dependent-type Lean proofs.}
\end{figure}

\section{What the two repositories change about your
part}\label{what-the-two-repositories-change-about-your-part}

\textbf{Decision:} integrate your work into \texttt{proof-graphs}; do
not build another parallel Lean graph extractor. Retain
\texttt{proofnet-ir} as the formally verified MLL reference layer.
Recast Boundary-Lens as a \textbf{typed, transaction-aware evidence and
proof-reuse interface}, evaluated against your friend's existing
structural-state and coupled-group search baselines.

This is a recommendation about research direction, not a judgment about
authorship or personal ability. The code is treated as your friend's
contribution because that is how you described it. The repositories
alone do not establish who contributed every idea, nor how publication
credit should be divided.

\subsection{The honest comparison}\label{the-honest-comparison}

At the inspected commits, your friend's work is ahead of the last
Boundary-Lens artifact in both implementation and experimental maturity.
Your earlier design remains useful, but much of its originally proposed
engineering is now overlapping work. The right response is to reuse that
foundation rather than defend an obsolete novelty claim. {[}R1--R8,
U1{]}

\begin{longtable}[]{@{}
  >{\raggedright\arraybackslash}p{(\columnwidth - 6\tabcolsep) * \real{0.2500}}
  >{\raggedright\arraybackslash}p{(\columnwidth - 6\tabcolsep) * \real{0.2500}}
  >{\raggedright\arraybackslash}p{(\columnwidth - 6\tabcolsep) * \real{0.2500}}
  >{\raggedright\arraybackslash}p{(\columnwidth - 6\tabcolsep) * \real{0.2500}}@{}}
\toprule\noalign{}
\begin{minipage}[b]{\linewidth}\raggedright
Dimension
\end{minipage} & \begin{minipage}[b]{\linewidth}\raggedright
\texttt{proofnet-ir}
\end{minipage} & \begin{minipage}[b]{\linewidth}\raggedright
\texttt{proof-graphs}
\end{minipage} & \begin{minipage}[b]{\linewidth}\raggedright
Your last available Boundary-Lens package
\end{minipage} \\
\midrule\noalign{}
\endhead
\bottomrule\noalign{}
\endlastfoot
Core object & Actual proof-net certificates for a restricted MLL
calculus & Extracted step-dependency graphs and live search graphs &
Goal-boundary snapshots and proposed typed graph updates \\
Trust basis & Lean-formalized checker/sequentialization results in that
calculus & Executed Lean traces, structural checks, final proof replay &
Included controlled Lean logs and custom frontier checks \\
Data scale & Formal-library and MLL experiment artifacts & 7,285
corrected step graphs, plus search artifacts & Five controlled theorems,
15 success rows and five failure rows in the inspected ZIP \\
Identity & Canonical identity for the specified MLL representation &
Lean-expression keys, renaming, and coupled groups & Raw goal transport;
schema and feature separation still incomplete \\
Search evidence & Registered restricted-logic experiments & Several
registered Lean search comparisons & No demonstrated incremental search
improvement \\
Best role in joint project & Logical reference and certified restricted
testbed & Shared implementation and empirical baseline & Proposed
evidence-quality, transport, and checked-reuse layer \\
\end{longtable}

The 7,285 graphs are not 7,285 fully annotated open-boundary traces. The
exported rows contain step kinds, source lines, edges, and graph
statistics; expression-rich state snapshots cannot be reconstructed from
that dataset alone. A sidecar enrichment must return to raw extraction
or re-elaborate a source subset. {[}R2, R7, R8, R17{]}

\subsection{What your friend has actually
contributed}\label{what-your-friend-has-actually-contributed}

\subsubsection{A. A genuine formal proof-net
library}\label{a.-a-genuine-formal-proof-net-library}

\texttt{proofnet-ir} establishes substantial guarantees for
\textbf{unit-free, cut-free multiplicative linear logic}: a correctness
checker, conversion from derivations, sequentialization of accepted
certificates, and canonical identity. Its rolling result includes a
stated quadratic symbolic-operation bound, not a general linear-time
guarantee or a bound on arbitrary Lean proof search. This is a
formal-methods contribution in its own right. {[}R1{]}

It is not simply a backend that can certify any Lean dependency graph.
Lean's dependent type theory and tactic meta-state are not MLL; a bridge
requires a separate definition and theorem. Shared terminology is not
such a bridge.

\subsubsection{B. An actual extraction and search
program}\label{b.-an-actual-extraction-and-search-program}

\texttt{proof-graphs} already traverses Lean InfoTrees, derives
dependency graphs, drives a live Lean REPL, compares state and goal
searches, normalizes identities structurally, and keeps
metavariable-coupled goals together. It also reconstructs proof scripts
and checks completed candidates from the statement. Those are
implementations, not only plans. {[}R5--R8{]}

Consequently, ``we use real Lean states,'' ``we preserve metavariable
identity,'' ``we group coupled goals,'' and ``we reuse solved goals''
are no longer distinctive next contributions for your branch.

\subsubsection{C. Useful negative
evidence}\label{c.-useful-negative-evidence}

The most valuable change is empirical. The repositories distinguish the
number of possible orderings of a completed proof graph from the order
redundancy actually encountered by a search policy. Default first-goal
scheduling can suppress the latter without eliminating the former.
{[}R2--R4{]}

In search-v0.8, the registered comparable set contains \textbf{158
tasks, with three draw sets per task}:

\begin{longtable}[]{@{}lrrrr@{}}
\toprule\noalign{}
Search & Set 0 & Set 1 & Set 2 & Total / 474 task-sets \\
\midrule\noalign{}
\endhead
\bottomrule\noalign{}
\endlastfoot
Whole states & 38 & 37 & 41 & 116 \\
Independent goals & 40 & 39 & 42 & 121 \\
Coupled goal groups & 40 & 39 & 43 & 122 \\
\end{longtable}

None of the three registered pairwise comparisons passed the
Holm-adjusted decision rule. These are neither 474 independent theorems
nor proof of equivalence between algorithms. They are evidence against
expecting a large automatic benefit from replacing state search by this
graph search in the tested setting. {[}R3{]}

The latest correction explicitly withdraws a simple runtime-cost
advantage: searches share model draws and run in a rotating order, so
later searches can inherit cached work. Both the pooled timing and the
``ran first'' subsets are inadequate as a clean isolated latency
comparison. {[}R3{]}

A further update matters in \texttt{proofnet-ir}: with the thinking
condition in representation-v0.3, 35 nets, 32 labelled derivations, and
30 positional derivations verify out of 90 tasks; no pairwise format
difference is significant after the adjustment. Without thinking the
counts are 11, 0, 0. An earlier headline that nets reliably outperform
derivations is too broad. {[}R9{]}

\subsection{Where your contribution still has
value}\label{where-your-contribution-still-has-value}

\subsubsection{Immediate, defensible value: independent semantic
audit}\label{immediate-defensible-value-independent-semantic-audit}

The existing handoff identifies unaudited failure combinators and
normalization concerns. A genuinely independent reconstruction and
adversarial Lean test suite is useful work, particularly because
structural invariants can pass when an extractor consistently assigns
the wrong interpretation. {[}R4{]}

Your Boundary-Lens perspective is valuable here: it asks what each edge
means, what state it belongs to, and what evidence permits reuse.
Convert that perspective into executable checks rather than another
layer of terminology.

\subsubsection{Medium-term value: a shared semantic
interface}\label{medium-term-value-a-shared-semantic-interface}

The offline extractor records goal identifiers and tactic-tree
structure, while the live search exporter records instantiated
expressions. A joint interface can connect these without treating them
as identical evidence: versioned goal capsules, transaction/branch
identities, assignment deltas, term/local provenance, and an explicit
record of incomplete dependency coverage. {[}R5--R7{]}

This is a proposed integration contribution. It is not yet implemented
in Lean in this deliverable, and it should not be advertised as a new
proof theory by naming it a lens.

\subsubsection{High-upside, uncertain value: transportable proof
fragments}\label{high-upside-uncertain-value-transportable-proof-fragments}

The new experiment should target a narrower question:

\begin{quote}
Can a boundary-aware index find and replay already checked proof
fragments in new compatible contexts, avoiding repeated tactic work
beyond structural-state deduplication and ordinary proof-term caching?
\end{quote}

Existing search already shares proofs for matching goal keys. The
possible extension is \textbf{typed context transport}: reuse when
irrelevant local assumptions differ, when required local evidence can be
mapped to a different context, or when the same subargument appears
inside different surrounding proofs. Start with closed fragments;
require Lean to check each transported application. Graph indexing must
beat a non-graph cache receiving the same fragment pool before
graph-specific value is claimed. {[}R6; proposed extension{]}

This has no guaranteed payoff. It may turn out that ordinary premise
retrieval or a typed cache captures almost everything. That outcome
would still determine which part of your representation is worth
keeping.

\subsection{Findings from this review}\label{findings-from-this-review}

A fresh Python audit of the available v0.7.1 ZIP found 15 success
records and five failure records, no typed edges in any success delta,
no explicit \texttt{rho\_before} in those rows, and 15/15 failures
against the shipped trace schema. That refines earlier conversation
counts: use the inspected artifact's five failures, not a remembered
six. No Lean execution was rerun here. {[}U1{]}

Reviewing \texttt{goal\_identity.py} also revealed a concrete
\textbf{testable hazard}: metavariable renaming uses regular expressions
over the complete serialized expression string, including quoted string
literals. Thus a literal containing a metavariable-looking token can be
altered as though it were a genuine metavariable. This package
reproduces that mechanism in small Python fixtures and shows how a
tagged AST avoids it. This is not evidence that a particular published
Lean proof was wrong or that the empirical conclusions are invalid. It
is a regression to verify in actual Lean. {[}R5{]}

Other bounded audit targets are an \texttt{h?} fallback for unexported
free variables, omitted implementation-detail locals, grouping by
directly visible term-metavariable mentions, and a printed carried-goal
check in the otherwise structural-key search. These deserve conservative
handling and tests. They are not a blanket claim that the entire engine
is unsound. Final proof replay protects theorem acceptance, but unsafe
merging can still discard solvable branches or distort measurements.
{[}R5, R6{]}

\subsection{What to retire from the old
pitch}\label{what-to-retire-from-the-old-pitch}

Retire ``factorial proof-order reduction implies factorial Lean search
savings.'' Retire ``state plus graph hashing safely prunes more than
faithful state hashing.'' Retire ``only the focused goal can change.''
Retire ``every successful tactic or recorded dependency edge is
kernel-certified.'' Keep the mathematical motivation, but distinguish
formal theorems, executed observations, derived annotations, and
speculative proposals.

Do not lead with a PFR or condensed-mathematics diagram whose states
were not extracted. Use those areas later as held-out semantic stress
tests, not as decorative evidence of scale.

\subsection{Proposed joint position}\label{proposed-joint-position}

A suitable joint research question is:

\begin{quote}
Which proof-graph abstractions remain useful once ordinary goal
scheduling and structural-state sharing have already removed the easy
redundancy, and what evidence is needed to reuse a proof fragment
safely?
\end{quote}

A suitable description of your part is:

\begin{quote}
I am developing the typed boundary and proof-transport layer that turns
extracted search graphs into auditable, reusable proof fragments, and
testing its incremental value against the project's existing state and
graph baselines.
\end{quote}

That is stronger than another general graph extractor because it
connects directly to the evidence your friend already collected. It is
also honest about what still has to be implemented and measured.

\section{Revised joint methodology: from graph shape to checked
reuse}\label{revised-joint-methodology-from-graph-shape-to-checked-reuse}

\textbf{Status:} proposed method and executable Python design tests. No
new Lean soundness theorem, neural-prover result, or frontier extraction
is claimed here. Mathematical propositions below have explicit
hypotheses; the reference implementation is a tagged-expression model,
not Lean's meta-state.

\subsection{Separate four objects}\label{separate-four-objects}

Use four related but distinct representations.

\begin{enumerate}
\def\labelenumi{\arabic{enumi}.}
\tightlist
\item
  \textbf{MLL certificate:} the object checked by \texttt{proofnet-ir}.
  Its correctness results apply only to its specified logic. {[}R1{]}
\item
  \textbf{Offline derivation graph:} the corrected \texttt{proof-graphs}
  step DAG. It is useful for structural analysis and supervised
  metadata. It is not an online state. {[}R7, R8{]}
\item
  \textbf{Live boundary/event graph:} accepted and attempted Lean
  transitions with source, state, and evidence identities. This is the
  proposed Boundary-Lens integration.
\item
  \textbf{Reuse index:} a retrieval structure for completed proof
  fragments. A match proposes an application; Lean decides whether that
  application is valid.
\end{enumerate}

This decomposition avoids making one graph serve simultaneously as a
proof certificate, a search state, a corpus row, and a heuristic index.

\subsection{Boundary and event
definitions}\label{boundary-and-event-definitions}

Let a live search configuration be

\[S=(E,O,Q,M,U,P),\]

where \(E\) is the environment and source-prefix availability, \(O\)
relevant options and scope, \(Q\) the ordered active-goal list, \(M\)
the reachable term-metavariable declarations and assignments, \(U\)
universe information, and \(P\) any relevant postponed or opaque
obligations. This is an interface contract; exporting all of Lean's
operational state is not assumed feasible. Missing fields must be
represented by coverage flags, not silently erased. Lean distinguishes
term goals, shared unknowns, and universe metavariables. {[}R14, R15{]}

A goal capsule stores its local telescope, target, available local
instances, identity, and content version. Metavariable types and
local-definition values participate in dependency closure. Raw
identifiers are for same-session provenance; model features use explicit
structural references. Do not replace an unexported free variable by a
common placeholder.

Define an observed boundary

\[\partial S=(Q,\operatorname{Reach}(Q;M,U,P),\operatorname{coverage}).\]

\texttt{Reach} follows references through goal targets, local
types/values, metavariable declarations/assignments, and universes.
Opaque constraints either join the affected component conservatively or
disable factorization. A successful export is not automatically a
complete export.

An event is

\[e_i=(b_i,S_i,t_i,o_i,S'_i,\eta_i),\]

with branch/transaction identifier \(b_i\), tactic syntax \(t_i\),
outcome \(o_i\), and evidence \(\eta_i\). Outcomes distinguish
unexecuted proposals, executed successes, errors, timeouts, and
unknowns. Disposition separately distinguishes committed, speculative,
and rolled-back events. A successful speculative alternative need not be
committed.

Let \(\rho_i\) map active Lean goal identifiers injectively into
persistent graph identities. Preserve this identity while incrementing a
\textbf{capsule version} when an assignment changes the instantiated
target or context. A removed active goal is \texttt{no\_longer\_active}
until assignment/reachability evidence determines whether it is solved,
replaced, or suspended.

The durable history is append-only:

\[H_{i+1}=H_i\mathbin{\|}[e_i].\]

The current graph is a materialized view:

\[G_{i+1}=\operatorname{Apply}(G_i,\Delta_i).\]

Node insertions can be monotone, but statuses, active ports, and content
versions are updates. Writing only \(G_{i+1}=G_i\cup\Delta_i\) loses
this distinction.

\subsection{Evidence hierarchy}\label{evidence-hierarchy}

Use independent evidence labels rather than a single \texttt{verified}
flag:

\begin{itemize}
\tightlist
\item
  \texttt{source\_observed}: syntax, arguments, source span;
\item
  \texttt{lean\_transition\_observed}: before/after state observed
  during an executed transition;
\item
  \texttt{assignment\_observed}: a particular assignment or expression
  reference was inspected;
\item
  \texttt{derived\_annotation}: an extractor rule inferred a relation;
\item
  \texttt{kernel\_replay\_passed}: a complete reconstructed declaration
  passed checking under a recorded trust policy;
\item
  \texttt{proposed}: no execution evidence.
\end{itemize}

An accepted tactic does not prove that a source-mentioned lemma was
used. \texttt{constructor} does not certify sibling independence. A root
proof with unresolved holes is not a complete theorem. Proof-term
enrichment reveals actual references but does not automatically
establish minimal dependencies or tactic commutation. {[}R14, R15{]}

\subsection{The hashing limitation: an elementary but decisive
result}\label{the-hashing-limitation-an-elementary-but-decisive-result}

\textbf{Proposition 1 (refinement cannot enlarge duplicate classes).}
Let \(N\) be any state key and \(H\) any graph key. Define
\(K(x)=(N(x),H(x))\). Then

\[K(x)=K(y)\Longrightarrow N(x)=N(y).\]

Therefore each \(K\)-equivalence class is a subset of an
\(N\)-equivalence class.

\textbf{Proof.} Equality of ordered pairs entails equality of their
first coordinates. \(\square\)

\textbf{Consequence.} A state-plus-graph exact key cannot, by itself,
merge more states than that same state key. It can retain \emph{more}
distinctions. Any advantage must come from a different decision:
scheduling, learned ranking, cross-context fragment reuse, or cheaper
comparison. Appending graph history to the duplicate key can even undo
useful transposition sharing.

This does not make graphs useless. It identifies which mechanism must be
measured. A faithful state may contain all logical information while a
graph remains a helpful computational representation. No
information-theoretic superiority is claimed.

\subsection{Conditional local independence, not tactic-name
independence}\label{conditional-local-independence-not-tactic-name-independence}

For an explicitly restricted tactic semantics, let \(R(t,S)\) and
\(W(t,S)\) denote the mutable cells read and written by a transition. If
both transitions are applicable, allocation can be renamed, they perform
no relevant external effects, and

\[W(t,S)\cap(R(u,S)\cup W(u,S))=\varnothing,\]
\[W(u,S)\cap(R(t,S)\cup W(t,S))=\varnothing,\]

with these footprints stable over the two executions, the usual
disjoint-update argument yields commuting results up to fresh-identifier
renaming.

This is a conditional frame argument, not a theorem about arbitrary Lean
tactics. A graph of visible goal references is an approximation to these
footprints. A shared witness can make a focused \texttt{exact} change
another goal. Arbitrary metaprograms may inspect additional state.
Consequently the production rule is: preserve observed identities,
record every observed effect, and refuse uncertified factorization when
coverage is insufficient. {[}R3, R5, R6, R14{]}

For \texttt{first}, \texttt{try}, and \texttt{repeat}, preserve nested
transaction ancestry. An InfoTree node in a failed alternative must not
be treated as a committed proof step solely because it is present in the
tree. The current repository's handoff explicitly leaves these
constructs as audit targets. {[}R4{]}

\subsection{The recommended new algorithm: closed-fragment
transport}\label{the-recommended-new-algorithm-closed-fragment-transport}

Begin with a completed local proof, not an arbitrary open graph. Write

\[E;\Gamma\vdash p:T,\]

where \(p\) contains no unresolved expression or universe metavariables.
Compute the dependency-closed local telescope \(\Gamma_p\): free locals
in \(p\) and \(T\), plus all locals required by their types and retained
local-definition values. Preserve order and instance/binder information.
Abstract to a portable term over \(\Gamma_p\); do not simply erase
unused-looking hypotheses from a printed state.

A candidate reuse at target \(E';\Gamma'\vdash T'\) supplies a typed
context substitution \(\theta\) and compatible environment. The proposed
application is \(\theta(p)\). The match is useful only if Lean accepts
it at \(T'\) under the current declaration prefix and trust policy.

\textbf{Proposition 2 (typed substitution, restricted setting).} In a
dependent type theory with the standard substitution rule, if
\(E;\Gamma\vdash p:T\) and \(\theta:\Gamma\to\Gamma'\) is a well-typed
context substitution, then

\[E;\Gamma'\vdash\theta(p):\theta(T).\]

\textbf{Proof sketch.} Induct through the typing derivation; the
variable case uses the substitution's typing assumptions, and binder
cases extend the substitution with a fresh bound variable. Local
definitions require compatibility with their values. If \(\theta(T)\) is
definitionally equal to \(T'\), the conversion rule gives the desired
judgment. \(\square\)

This is standard metatheory, not an original theorem claimed by the
project. It motivates a concrete checker: build the candidate in Lean
rather than trusting Python's matcher. Different environments require
checking the used declarations, universe parameters, and permitted
axioms; equal spelling of a constant is not sufficient.

\subsubsection{Online procedure}\label{online-procedure}

\begin{Shaded}
\begin{Highlighting}[]
\NormalTok{For each actually solved local obligation:}
\NormalTok{  inspect the resulting term in its saved Lean context;}
\NormalTok{  reject unresolved holes, unexported locals, or unsupported obligations;}
\NormalTok{  build its dependency{-}closed telescope and portable abstraction;}
\NormalTok{  store term reference, interface, provenance, and replay evidence.}

\NormalTok{Before requesting another costly tactic batch:}
\NormalTok{  query the ordinary typed cache and/or the boundary{-}graph index;}
\NormalTok{  propose context substitutions using only the available prefix;}
\NormalTok{  instantiate a candidate in a cloned current Lean state;}
\NormalTok{  accept a reuse only if Lean accepts the application;}
\NormalTok{  commit exactly the accepted branch and record all state effects;}
\NormalTok{  otherwise restore state and continue the unchanged baseline search.}
\end{Highlighting}
\end{Shaded}

The index may be approximate. False matches are performance failures,
not licenses to assert a proof. Keep final full-declaration replay.
Typechecking the reused term protects theorem correctness; maintaining
baseline fallbacks protects coverage but does not by itself guarantee
equal success under a finite budget.

\begin{figure}
\centering
\includegraphics[width=0.68\textwidth,height=\textheight]{../figures/reuse_stages.pdf}
\caption{Stages of a proposed checked reuse. Retrieval and context
matching remain proposals until Lean accepts the application.}
\end{figure}

\subsection{Three instructive
examples}\label{three-instructive-examples}

\subsubsection{A. Existing sibling transport is a regression, not a new
speed
result}\label{a.-existing-sibling-transport-is-a-regression-not-a-new-speed-result}

From \(h:p\land q\), use \texttt{constructor}, \texttt{exact\ h.right},
\texttt{exact\ h.left} to prove \(q\land p\). The second tactic must
preserve the identity of the remaining \(p\) goal. Both the user
artifact and the current research line already target this case. It
tests extraction, not an incremental graph-search advantage.

\subsubsection{B. Context-sensitive reuse: the actual new
hypothesis}\label{b.-context-sensitive-reuse-the-actual-new-hypothesis}

Suppose a completed local term proves \(f\ a=g\ a\) from \(h:f=g\). In
another branch, the context includes extra unrelated assumptions and
differently named \(f,g,a,h\). The whole context key can differ even
though the fragment
\texttt{congrArg\ (fun\ k\ =\textgreater{}\ k\ a)\ h} transports. A
dependency-indexed cache can propose the fragment using only its
required telescope; Lean checks the application.

Graph stages are: an open equality port; a retrieved fragment with
unbound interface parameters; a proposed typed substitution; a
successful local application; and the unchanged external frontier. This
could avoid re-searching the local argument. However an ordinary
proof-term cache could do exactly the same, so it is a mandatory
baseline. This is a design example, not a measured result.

\subsubsection{C. Coupled witnesses: why local successes cannot be
combined
freely}\label{c.-coupled-witnesses-why-local-successes-cannot-be-combined-freely}

For an open construction of \(\exists n:\mathbb N,\ n=0\land n=1\), a
shared witness cell \(w\) links the residual goals \(w=0\) and \(w=1\).
One isolated branch can succeed by setting \(w:=0\), another by setting
\(w:=1\). Their union is inconsistent. A reuse packet for an open
fragment must include assignment constraints and jointly check
composition; a cache that stores only ``this goal succeeded'' is
insufficient. This is why closed-fragment reuse precedes open-fragment
reuse.

\begin{figure}
\centering
\includegraphics[width=0.85\textwidth,height=\textheight]{../figures/shared_witness.pdf}
\caption{Why a shared witness must remain in the interface of an open
fragment. This is a logical counterexample to naive composition, not an
executed Lean trace.}
\end{figure}

\subsection{Interface and integration}\label{interface-and-integration}

Keep the frozen v0.2 graph dataset unchanged. Add
\texttt{joint-sidecar/0.1.0} records keyed by source commit, declaration
locator, extraction session, InfoTree/event path, and content hash. The
existing exporter does not provide rich boundary state, so its imported
sidecars truthfully contain \texttt{state\_before:\ null},
\texttt{state\_after:\ null}, and
\texttt{boundary\_coverage:\ not\_exported}. {[}R7, R8, R17{]}

Use a tagged JSON expression codec rather than renaming over serialized
strings. Preserve literals and alias structure; handle missing
declarations as unknown. Use full digests for indexing and retain
canonical bytes for collision resolution. Neither a digest nor a
transport map is a proof of Lean-state equivalence.

The reference code in this package demonstrates those representation
properties in a finite model. It deliberately rejects incomplete fixture
coverage. It does not supply a Lean exporter, a complete operational
boundary, a proof-term matcher, or kernel replay.

\subsection{Scope of novelty}\label{scope-of-novelty}

Existing work supplies the MLL theory, hypergraph search, Lean
interfaces, factorized state reuse, and expression graphs. {[}R1,
R10--R16{]} This proposal adapts those ideas into an auditable
composition of state snapshots, event effects, portable proof fragments,
and checked reuse. The original deliverable here is the integration
design and its adversarial regression model, not a claim to have
invented substitution, proof graphs, metavariable coupling, or
memoization.

Tactician's Web already connects Coq definitions, proof terms, tactics,
and states in a semantic graph with an interactive interface; Graph2Tac
and Tactician's online learning show that reuse of nearby formal
knowledge is also an established research theme. {[}R18, R19{]} These
precedents further rule out claiming generic semantic proof graphs or
online reuse as new. The proposed contribution must be its precise
Lean-side contract, independent evidence, and measured incremental
behavior.

Only experiments can establish whether the joint representation adds
useful value. If a plain typed cache matches graph-indexed reuse, keep
the cache and report that the graph-specific hypothesis failed.

\section{Experiment and implementation
plan}\label{experiment-and-implementation-plan}

\textbf{Registration status: DRAFT.} This document is not a
preregistration. Freeze code, corpus, hypotheses, exclusions, costs, and
analysis in a new commit before collecting new outcome data. Do not edit
existing registered experiments or reinterpret their decisions. {[}R3,
R4{]}

\subsection{Work package 1: reproduce and audit the actual
baseline}\label{work-package-1-reproduce-and-audit-the-actual-baseline}

Work in a local branch of \texttt{proof-graphs}, pinned initially at
\texttt{8010d137af5af5040fcbcd9e98f910b01162263e}. Keep its actual
dependency lock; do not substitute the newer \texttt{proofnet-ir} main
for a pinned dependency. Compare that pin with the separately inspected
ProofNet-IR head and document the distinction.

Run the documented build and committed-artifact checks. Distinguish
recomputing summaries from running new inference. Check a small set of
fresh Lean fixtures before a broad export. Select fixtures for nested
\texttt{have}, \texttt{calc}, \texttt{simpa\ ...\ using},
\texttt{first}, \texttt{try}, \texttt{repeat}, shared witnesses,
metavariable-dependent local types, and universe sharing. A human should
independently reconstruct at least the failure-combinator cases before
viewing extractor output. {[}R4{]}

Acceptance: recorded commands and logs; untouched frozen outputs; a
reproducible failing test for each confirmed defect, or an explicit
refutation of the suspected defect. No mass source transformation or
model training is required.

\subsection{Work package 2: a fail-closed typed boundary
sidecar}\label{work-package-2-a-fail-closed-typed-boundary-sidecar}

Extend the existing Lean exporter and REPL rather than introducing
another session driver. Required information is the active goal list,
persistent raw IDs, versioned capsules, tagged expressions and levels,
local telescope, reachable metavariable declarations/assignments,
environment/prefix fingerprint, branch identity, and dependency-coverage
flags.

When complete coverage is unavailable, keep the entire affected
state/group and avoid hard equality merges based on lossy printed
fallback. Lack of factoring is an acceptable conservative result.
Querying the same snapshot twice must not change observable proof
behavior.

Acceptance: all adversarial fixtures pass; no literal modification or
placeholder identity merges; same raw goal may change its capsule
version; failures do not mutate the committed graph. Every exported row
validates against one current schema. Recorded gold evidence is attached
only to actual full replay, not to a tactic name.

\subsection{Work package 3: closed-fragment cache with a graph
ablation}\label{work-package-3-closed-fragment-cache-with-a-graph-ablation}

Use only completed, hole-free local terms. Start with same-environment
reuse under typed renaming and weakening. Do not start with arbitrary
open-subproof unification or a learned GNN.

\textbf{Arms:}

\begin{longtable}[]{@{}
  >{\raggedright\arraybackslash}p{(\columnwidth - 2\tabcolsep) * \real{0.5000}}
  >{\raggedright\arraybackslash}p{(\columnwidth - 2\tabcolsep) * \real{0.5000}}@{}}
\toprule\noalign{}
\begin{minipage}[b]{\linewidth}\raggedright
Arm
\end{minipage} & \begin{minipage}[b]{\linewidth}\raggedright
Purpose
\end{minipage} \\
\midrule\noalign{}
\endhead
\bottomrule\noalign{}
\endlastfoot
A0: faithful whole-state search & Existing baseline, same correctness
repairs as other arms \\
A1: faithful coupled-group search & Existing graph baseline \\
A2: A0 plus ordinary typed closed-term cache & Does reuse help without
graph-specific indexing? \\
A3: A0 plus boundary/graph-indexed cache & Does your representation
improve retrieval or net cost? \\
\end{longtable}

A2 and A3 receive the same proof-fragment inventory, trust policy,
verification budget, and candidate proposer. A3 cannot quietly receive
more theorems, future dependencies, or a stronger model. Repair keying
in all arms; do not credit fixing an old baseline bug as a
graph-specific learning benefit.

Primary mechanism metrics: eligible repeated fragments, retrieval
attempts, successful Lean-checked reuse, rejected transports, avoided
model requests, avoided tactic executions, export/match/recheck costs.
Primary task metric: theorem success under an explicitly fixed resource
budget. Report all denominators and setup failures.

A graph feature could improve retrieval order without supporting any
extra exact merges. Make that the experimental distinction rather than
comparing a weak string cache to a rich graph cache.

\subsection{Work package 4: held-out mathematical
pilot}\label{work-package-4-held-out-mathematical-pilot}

First use 20--50 small and medium proof fragments with genuine repeated
subarguments and adversarial negatives; choose the final count before
the run. After extraction and transport are stable, select a separate
held-out module set in analysis, algebra, or category theory. Curate by
syntactic constructs and available build resources, not by whether the
new algorithm wins.

For condensed mathematics, use an actually buildable, pinned Lean
project and identify its Lean generation. The historical Liquid Tensor
corpus and a modern Lean 4 Mathlib development are not interchangeable.
No theorem from that area was extracted in this review. The milestone is
one correct traced module fragment, not a diagram labelled
frontier-scale.

\subsection{Controls that matter}\label{controls-that-matter}

\textbf{Leakage.} Split by module/theorem family and dependency
availability. Exclude the target declaration itself and later
declarations from the proving environment. Store full proofs for audit
separately from policy inputs. Do not compute input graph features from
future proof steps.

\textbf{Common draws.} For controlled search comparison, prerecord or
deterministically index candidate draws by exact request content, seed,
and occurrence. Report marginal candidate costs rather than which arm
happened to populate the shared cache. A policy conditioned on a
different graph prompt cannot be called the same draw experiment without
describing the changed request.

\textbf{Runtime.} Run isolated fresh-process cold-cache trials,
randomized by arm order, under an explicit warmup policy. Include
extraction, serialization, canonicalization, matching, replay, and model
time. The corrected search-v0.8 cost discussion makes this a required
control. {[}R3{]}

\textbf{Statistics.} A theorem repeated under seeds is a clustered
observation. Use paired per-task comparisons and module-aware
uncertainty estimates. Predefine multiple-comparison handling.
Nonsignificance is not equivalence; a claim of non-inferiority needs a
chosen margin and an appropriate design.

\subsection{Go / no-go gates}\label{go-no-go-gates}

Proceed to a trained graph model only after there is evidence of a
useful signal beyond a plain cache. If eligible reuse is rare, change
corpus or conclude limited applicability before optimizing graph
encoders. If replay overhead exceeds avoided work, report the cost
failure. If A2 and A3 perform alike, remove graph-specific claims and
retain the useful semantic tooling.

The immediate decision need not depend on a positive ML result: a
carefully audited extractor and falsifiable account of where abstraction
fails are legitimate outputs, though their publication fit depends on
the resulting evidence and venue.

\section{Contribution and collaboration
proposal}\label{contribution-and-collaboration-proposal}

\subsection{Your part, stated without inflating
it}\label{your-part-stated-without-inflating-it}

Your last available implementation is smaller and less mature than your
friend's current repositories. Its strengths are the original problem
framing, the focus on partial boundaries rather than finished proofs
alone, and the insistence on extracting faithful intermediate objects.
Those strengths become a distinct contribution only when they are turned
into new checks, semantics, or demonstrated reuse---not when renamed as
another graph abstraction. {[}R1--R8, U1{]}

This assessment compares artifacts. It cannot reconstruct the history of
your joint conversations, the origin of shared ideas, or individual
authorship. Agree on contribution credit directly and keep a written
record of work performed.

\subsection{Proposed ownership, subject to your
agreement}\label{proposed-ownership-subject-to-your-agreement}

\begin{longtable}[]{@{}
  >{\raggedright\arraybackslash}p{(\columnwidth - 4\tabcolsep) * \real{0.3333}}
  >{\raggedright\arraybackslash}p{(\columnwidth - 4\tabcolsep) * \real{0.3333}}
  >{\raggedright\arraybackslash}p{(\columnwidth - 4\tabcolsep) * \real{0.3333}}@{}}
\toprule\noalign{}
\begin{minipage}[b]{\linewidth}\raggedright
Workstream
\end{minipage} & \begin{minipage}[b]{\linewidth}\raggedright
Suggested lead
\end{minipage} & \begin{minipage}[b]{\linewidth}\raggedright
Concrete output
\end{minipage} \\
\midrule\noalign{}
\endhead
\bottomrule\noalign{}
\endlastfoot
Restricted MLL theory and certified IR & Friend / existing maintainer &
Stable theorem interface; no unsupported transfer to Lean graphs \\
Existing extraction, search, and frozen results & Friend / existing
maintainer & Baseline behavior, reproducibility support, review of
changes \\
Independent semantic audit & You & Hand reconstructions, adversarial
Lean fixtures, reproducible bug reports \\
Typed boundary / event interface & You, jointly reviewed & Versioned
schema, exporter sidecar, provenance, transaction checks \\
Checked fragment transport & You, jointly reviewed & Closed-fragment
cache, context mapping, replay evidence \\
Experimental design and interpretation & Both & Frozen protocol, strong
baselines, negative results retained \\
Paper and release & Both & Separate formal claims, engineering
guarantees, empirical conclusions \\
\end{longtable}

Do not divide the project as ``he does code; I do ideas.'' Own a module
with measurable acceptance criteria. Conversely, a correctness audit is
not inferior to new features: it determines whether the features support
the intended scientific claims.

\subsection{Proposed sequence of pull
requests}\label{proposed-sequence-of-pull-requests}

\textbf{PR 1: audit only.} New adversarial fixtures and source-pinned
review. No baseline behavior change. Establish which suspicious cases
are real defects.

\textbf{PR 2: typed boundary sidecar.} Additive schema and serializer,
identity-version separation, conservative unknown handling. Keep old
datasets and frozen experiments intact.

\textbf{PR 3: checked closed-fragment cache.} Baseline cache and
graph-index variant using the same inventory. No learned model yet.

\textbf{PR 4: registered experiment and report.} Collected only after
the protocol commit. Explain null or negative results as clearly as
positive ones.

No automatic push to main. Work on a local feature branch, request
review, and resolve concurrent-edit ownership. The final ZIP is a
planning/reference package, not a modification of either repository.

\subsection{Suggested message to your
friend}\label{suggested-message-to-your-friend}

``I read the newer results and think I should stop duplicating the
extractor and goal-grouping work you already have. I'd like to own the
semantic audit and typed boundary/transport layer, then test whether
checked reuse of local proof fragments adds anything beyond your
structural-state and group baselines. I will start with the
failure-combinator and identity tests, keep your frozen experiments
intact, and include a plain proof-term cache so we can isolate whether
the graph really helps. Let's agree on that scope and review the
interfaces together before implementation.''

\section{What was and was not
executed}\label{what-was-and-was-not-executed}

\subsection{Actually executed here}\label{actually-executed-here}

The supplied user v0.7.1 archive was read programmatically and checked
against its shipped JSON Schema.
\texttt{evidence/USER\_V071\_AUDIT.json} records the counts and mismatch
examples. This is a fresh artifact audit, not a fresh Lean run.

The package's Python reference tests were run with
\texttt{python\ -m\ unittest\ discover\ -s\ tests\ -v}. They cover typed
renaming, literal preservation, metavariable aliasing, shared local
identity across capsules, transitive and universe dependencies in a
finite model, frontier persistence, and schema guards.
\texttt{evidence/REFERENCE\_TEST\_LOG.txt} is the actual output.

The literal regression reproduces the specific regex-renumbering
mechanism observed in the pinned \texttt{scripts/goal\_identity.py}
{[}R5{]}. It does not import and execute the entire remote repository,
and its synthetic inputs are not evidence of an actual erroneous Mathlib
proof.

\subsection{Not executed here}\label{not-executed-here}

No local \texttt{lake\ build}, Lean proof elaboration, friend-repository
end-to-end reproduction, full 7,285-row dataset conversion, new model
call, or frontier theorem trace was performed. Lean/Lake was
unavailable; container GitHub cloning failed at DNS resolution, while
selected files were successfully read through the connected GitHub tool.
No repository was modified remotely.

\texttt{examples/BoundaryRegression.lean} is uncompiled fixture source.
The JSON examples are explicitly synthetic design fixtures. Their IDs
name fixtures, not invented Lean metavariables. The offline adapter was
tested on a synthetic row matching the inspected format, not the full
downloaded dataset.

\subsection{Interpretation limits}\label{interpretation-limits}

The Python representation is not a Lean typechecker and does not decide
operational state equivalence. It deliberately rejects missing declared
dependencies, but its finite-model closure does not establish complete
coverage of Lean's postponed constraints, plugins, transparency
behavior, elaborator state, or external effects. A matching fingerprint
licenses no hard pruning. The sidecar schema enforces that limitation.

A passed JSON schema checks structure, not provenance truth. A claimed
log reference still requires checking its contents. A passed local term
application is not a full-declaration replay. A full checked theorem
does not prove that every extractor annotation is correct or that a
pruning rule preserves search completeness.

\section{Source register}\label{source-register}

Reviewed 30 September 2026. GitHub links are pinned to the exact
inspected commit. The Lean reference is a moving documentation URL;
actual integration must inspect the pinned Lean source. R-numbers below
are used throughout the reports. U1 is a local artifact, not a public
publication.

\textbf{{[}R1{]} ProofNet-IR: current status and restricted MLL
guarantees.}
\url{https://github.com/fushanbobfan/proofnet-ir/blob/261f2db26aac842534f4306da0f45d8c5e9b1c05/docs/current-status.md}

\textbf{{[}R2{]} proof-graphs: current scope and dataset.}
\url{https://github.com/fushanbobfan/proof-graphs/blob/8010d137af5af5040fcbcd9e98f910b01162263e/README.md}

\textbf{{[}R3{]} proof-graphs: states, goals, and coupled groups,
search-v0.8.}
\url{https://github.com/fushanbobfan/proof-graphs/blob/8010d137af5af5040fcbcd9e98f910b01162263e/experiments/search-v0.8/README.md}

\textbf{{[}R4{]} proof-graphs: handoff, verification boundaries, and
audit priorities.}
\url{https://github.com/fushanbobfan/proof-graphs/blob/8010d137af5af5040fcbcd9e98f910b01162263e/HANDOFF.md}

\textbf{{[}R5{]} proof-graphs: structural goal identity implementation.}
\url{https://github.com/fushanbobfan/proof-graphs/blob/8010d137af5af5040fcbcd9e98f910b01162263e/scripts/goal_identity.py}

\textbf{{[}R6{]} proof-graphs: structural-key whole-state, goal, and
group searches.}
\url{https://github.com/fushanbobfan/proof-graphs/blob/8010d137af5af5040fcbcd9e98f910b01162263e/scripts/search_faithful.py}

\textbf{{[}R7{]} proof-graphs: Lean InfoTree extractor.}
\url{https://github.com/fushanbobfan/proof-graphs/blob/8010d137af5af5040fcbcd9e98f910b01162263e/ProofGraphs/Extract.lean}

\textbf{{[}R8{]} proof-graphs: corrected graph derivation and dataset
generation.}
\url{https://github.com/fushanbobfan/proof-graphs/blob/8010d137af5af5040fcbcd9e98f910b01162263e/scripts/run_recount.py}

\textbf{{[}R9{]} ProofNet-IR: representation study with and without
thinking, v0.3.}
\url{https://github.com/fushanbobfan/proofnet-ir/blob/261f2db26aac842534f4306da0f45d8c5e9b1c05/experiments/representation-v0.3/README.md}

\textbf{{[}R10{]} Kripner, Sustr, and Straka (2025), LeanTree:
Accelerating White-Box Proof Search with Factorized States in Lean 4.}
\url{https://arxiv.org/abs/2507.14722v1}

\textbf{{[}R11{]} Lample et al.~(2022), HyperTree Proof Search for
Neural Theorem Proving.} \url{https://arxiv.org/abs/2205.11491v1}

\textbf{{[}R12{]} Aniva, Sun, Miranda, Barrett, and Koyejo (2025
revision), Pantograph.} \url{https://arxiv.org/abs/2410.16429v2}

\textbf{{[}R13{]} Aniva, Oikawa, Dill, and Barrett (2026), Nazrin: An
Atomic Neural Proof Automation Tactic in Lean 4.}
\url{https://arxiv.org/abs/2602.18767v3}

\textbf{{[}R14{]} Lean language reference: reading proof states.}
\url{https://lean-lang.org/doc/reference/latest/Tactic-Proofs/Reading-Proof-States/}

\textbf{{[}R15{]} Lean language reference: tactic proofs.}
\url{https://lean-lang.org/doc/reference/latest/Tactic-Proofs/}

\textbf{{[}R16{]} Yang et al.~(2023), LeanDojo: Theorem Proving with
Retrieval-Augmented Language Models.}
\url{https://arxiv.org/abs/2306.15626}

\textbf{{[}R17{]} proof-graphs: original dataset row format; use R8 for
corrected v0.2 generation.}
\url{https://github.com/fushanbobfan/proof-graphs/blob/8010d137af5af5040fcbcd9e98f910b01162263e/scripts/export_graphs.py}

\textbf{{[}U1{]} User-supplied Boundary-Lens v0.7.1 archive; fresh
Python audit in this package.} ../evidence/USER\_V071\_AUDIT.json

The repository reports are first-party evidence, not independently
replicated results in this review. Selected code paths were inspected
through the connected GitHub reader. Container cloning was blocked by
DNS and Lean/Lake was unavailable. The review therefore does not certify
either full repository build or all experimental conclusions.

\textbf{{[}R18{]} Blaauwbroek et al.~(2024), Graph2Tac: Online
Representation Learning of Formal Math Concepts, ICML / PMLR 235.}
\url{https://proceedings.mlr.press/v235/blaauwbroek24a.html}

\textbf{{[}R19{]} Blaauwbroek (2024), The Tactician's Web of Large-Scale
Formal Knowledge.} \url{https://arxiv.org/abs/2401.02950v2}

\end{document}
